\section{Discussion}
\label{sec:discussion}

\paragraph{What to send a model.} Extract the text. A PDF, DOCX, PPTX or
XLSX read through a standard loader costs what plain text costs and
answers as well. The loader takes milliseconds. Raw office XML costs
4--8$\times$ the tokens and, at any fixed budget, most of the answers. JSON
and HTML are worth stripping when context is tight: they lose 20--30
points at the plain-text budget. With a large context and a capable
model, format matters little to accuracy. It still matters to time,
memory and money, which all scale with tokens, and on a local GPU
possibly much faster than linearly.

\paragraph{What to ask a model for.} Ask for Markdown, then convert. A PDF
through Markdown cost as much as the Markdown itself. A PDF through LaTeX
cost \TaxPdfViaTex$\times$ for the 8B model and was never produced by the
two smaller models. LaTeX's failures were about escaping, not layout. Office files
are the expensive end. The cost is not in the file but in the code that
makes it, which the model has to get right against a library API, without
seeing the result. A slide deck was the worst case on every measure. An
agent that has to deliver a .docx or .pptx is better served by writing
Markdown and converting it with a fixed, tested tool (pandoc, a template)
than by writing python-docx or python-pptx code each time.

\paragraph{Raw XML is a surprise as an output.} The 8B model never failed
to produce a valid \texttt{document.xml}, while its python-docx code
needed repair in most cases. For a capable model, the format with no API
to misremember was the reliable one. It was not the faithful one, though:
\RecallDocxXml\,\% of the facts against \RecallDocxCode\,\% for code.

\paragraph{Limitations.}
\begin{itemize}
\item \textbf{Small, local models.} The largest model has 8B parameters,
  in 4-bit. Frontier models would fail less, and their tax would shrink
  towards the token ratio. Token counts do not
  depend on model quality. The write results are a lower bound on
  reliability.
\item \textbf{Synthetic documents.} The reports are short (about
  \TxtTokens{} tokens of plain text) and regular. Real office files carry
  styles, revision marks and embedded objects that make raw XML heavier
  than ours, and longer documents would reach the memory cliff sooner.
\item \textbf{Few documents for writing.} Writing used
  \NWriteDocsQwenTwoFiveOneFiveb{} documents for the 1.5B model,
  \NWriteDocsQwenThreeEightb{} for the 8B model (the sixth ran out of
  memory) and \NWriteDocsPhiThreeFiveMini{} for Phi-3.5-mini, which we
  stopped after 7 hours: past the memory cliff, one 8-section document
  would have taken most of a day. Validity rates are coarse at that size, and the tax ratios
  carry the variance of a few slow failures.
\item \textbf{One stack.} The memory cliff belongs to one GPU, one
  attention implementation and one driver. Its position will move with
  all three. Its existence will not, as long as markup adds tokens.
\item \textbf{Fidelity is string recall.} It checks that names and numbers
  survive, not layout, styling or whether a slide reads well.
\end{itemize}
